\section{Simulated antigens}\label{sec:simulated-antigens}

One section per antigen we simulate: what it is, the structural handle we use, and its
alanine-scan labels. Each structure is coloured per residue by \ddG\ on a diverging scale
with \textbf{white pinned at $0$}: red $=$ positive (hotspot / likely epitope), blue $=$
negative, grey $=$ residue not scanned. One panel per antibody; a single global colour scale
($\pm$max$|\ddG|$ over the whole benchmark, white at $0$) is shared by every figure, so colours
are directly comparable across antigens. Vertical colourbar at right. \ddG\ is the change in
binding free energy, $K_d$/$K_i$-derived; positive weakens binding.

\subsection{Lysozyme (hen egg-white, HEL)}
Hen egg-white lysozyme is a 129-residue glycoside hydrolase and the canonical model antigen of
structural immunology: multiple high-resolution antibody complexes and alanine scans exist
(HyHEL-10, HyHEL-63, D1.3). It is compact and conformationally rigid, cross-linked by four
disulfides. We simulate the apo enzyme from \texttt{1AKI}. The labels here are the HyHEL-63 scan
(\texttt{1DQJ}, Table~\ref{tab:ddg-lysozyme})~\cite{li2000}; the v2 benchmark adds D1.3
(\texttt{1VFB})~\cite{bhat1994} and HyHEL-10 (\texttt{3HFM})~\cite{padlan1989} scans for the same
antigen. These $\ddG$ values are drawn from the AB-Bind database~\cite{sirin2016}.
\begin{figure}[H]\centering
\begin{minipage}[c]{0.27\linewidth}\centering\figorbox{lysozyme_1DQJ_ddg.png}{1.0}\\[-2pt]{\small HyHEL-63 (1DQJ)}\end{minipage}\hfill
\begin{minipage}[c]{0.27\linewidth}\centering\figorbox{lysozyme_1VFB_ddg.png}{1.0}\\[-2pt]{\small D1.3 (1VFB)}\end{minipage}\hfill
\begin{minipage}[c]{0.27\linewidth}\centering\figorbox{lysozyme_3HFM_ddg.png}{1.0}\\[-2pt]{\small HyHEL-10 (3HFM)}\end{minipage}\hfill
\begin{minipage}[c]{0.13\linewidth}\centering\figorbox{ddgbar.png}{1.0}\end{minipage}
\caption{HEL coloured per residue by alanine-scan \ddG\ for its three antibodies (shared scale, white $=0$, red $=$ hotspot, blue $<0$; grey $=$ not scanned). Distinct but overlapping epitope cores across HyHEL-63, D1.3 and HyHEL-10.}\end{figure}
% AUTO-GENERATED by analysis/make_label_tables.py -- do not edit by hand.
\begin{table}[htbp]\centering\small
\caption{\textbf{Lysozyme} alanine-scan \ddG\ (kcal/mol; SPR, $K_d$-derived via $\ddG=RT\ln(K^{\mathrm{mut}}/K^{\mathrm{wt}})$, $T=310$\,K; start structure: apo 1AKI). Positive \ddG\ weakens binding.}
\label{tab:ddg-lysozyme}
\begin{tabular}{lr}
\toprule
Residue & HyHEL-63 (1DQJ) \\
\midrule
Y20 & +3.30 \\
R21 & +1.30 \\
W62 & +0.80 \\
W63 & +1.30 \\
L75 & +1.50 \\
T89 & +0.80 \\
N93 & +0.60 \\
K96 & +6.10 \\
K97 & +3.50 \\
S100 & +0.80 \\
D101 & +1.30 \\
\bottomrule
\end{tabular}
\end{table}


\subsection{HPr}
The histidine-containing phosphocarrier protein (HPr) is a small ($\sim$9\,kDa, 85-residue)
cytoplasmic protein of the bacterial phosphoenolpyruvate:sugar phosphotransferase system,
shuttling a phosphoryl group between enzyme~I and the sugar-specific EIIA components. The
\textit{E.\ coli} protein is a well-folded open-faced $\beta$-sandwich with \emph{no cysteines},
making it an exceptionally clean dynamics test case. The Jel42 epitope clusters in the C-terminal
half (Table~\ref{tab:ddg-hpr})~\cite{prasad1998}; $\ddG$ values via AB-Bind~\cite{sirin2016}. We
simulate the apo protein from \texttt{2JEL} (chain P).
\begin{figure}[H]\centering
\begin{minipage}[c]{0.6\linewidth}\centering\figorbox{hpr_2JEL_ddg.png}{0.95}\\[-2pt]{\small Jel42 (2JEL)}\end{minipage}\hfill
\begin{minipage}[c]{0.13\linewidth}\centering\figorbox{ddgbar.png}{1.0}\end{minipage}
\caption{HPr coloured per residue by Jel42 alanine-scan \ddG\ (white $=0$, red $=$ hotspot; grey $=$ not scanned). A single dominant hotspot (E70).}\end{figure}
% AUTO-GENERATED by analysis/make_label_tables.py -- do not edit by hand.
\begin{table}[htbp]\centering\small
\caption{\textbf{HPr} alanine-scan \ddG\ (kcal/mol; competition, $K_d$-derived via $\ddG=RT\ln(K^{\mathrm{mut}}/K^{\mathrm{wt}})$, $T=310$\,K; start structure: strip 2JEL chain P). Positive \ddG\ weakens binding.}
\label{tab:ddg-hpr}
\begin{tabular}{lr}
\toprule
Residue & Jel42 antibody (2JEL) \\
\midrule
T62 & +0.00 \\
E68 & +0.41 \\
E70 & +2.75 \\
H76 & -0.41 \\
E83 & +0.00 \\
E85 & +0.00 \\
\bottomrule
\end{tabular}
\end{table}


\subsection{VEGF-A}
Vascular endothelial growth factor A is a secreted pro-angiogenic cytokine and a central drug
target (bevacizumab, ranibizumab). Its receptor-binding domain is an antiparallel,
disulfide-linked homodimer with a cystine-knot growth-factor fold; the two protomers are
covalently joined, so the biological unit must be simulated as the dimer. The benchmark pairs two
antibodies against the same receptor-binding face: fab-12 (\texttt{1BJ1}, the bevacizumab
precursor)~\cite{muller1998} and the affinity-matured Y0317 (\texttt{1CZ8},
ranibizumab)~\cite{chen1999} --- a useful same-antigen/different-affinity comparison
(Table~\ref{tab:ddg-vegf}); $\ddG$ values via AB-Bind~\cite{sirin2016}. We simulate the apo dimer
from \texttt{1BJ1} (chains V,\,W).
\begin{figure}[H]\centering
\begin{minipage}[c]{0.37\linewidth}\centering\figorbox{vegf_1BJ1_ddg.png}{1.0}\\[-2pt]{\small fab-12 (1BJ1)}\end{minipage}\hfill
\begin{minipage}[c]{0.37\linewidth}\centering\figorbox{vegf_1CZ8_ddg.png}{1.0}\\[-2pt]{\small Y0317 (1CZ8)}\end{minipage}\hfill
\begin{minipage}[c]{0.13\linewidth}\centering\figorbox{ddgbar.png}{1.0}\end{minipage}
\caption{VEGF-A homodimer coloured per residue by \ddG\ for each antibody (shared scale, white $=0$; grey $=$ not scanned). Both protomers carry the same labelled residues; red marks the hotspots.}\end{figure}
% AUTO-GENERATED by analysis/make_label_tables.py -- do not edit by hand.
\begin{table}[htbp]\centering\small
\caption{\textbf{VEGF} alanine-scan \ddG\ (kcal/mol; SPR, $K_d$-derived via $\ddG=RT\ln(K^{\mathrm{mut}}/K^{\mathrm{wt}})$, $T=310$\,K; start structure: strip 1BJ1 chains V,W (dimer)). Positive \ddG\ weakens binding.}
\label{tab:ddg-vegf}
\begin{tabular}{lrr}
\toprule
Residue & fab-12 (1BJ1) & Y0317 (1CZ8) \\
\midrule
F17 & +0.00 & +0.00 \\
Y21 & +0.00 & +0.00 \\
Y45 & +0.82 & +1.94 \\
K48 & +0.41 & +0.00 \\
Q79 & +0.00 & +0.65 \\
I80 & +0.82 & +0.95 \\
M81 & +3.69 & +4.10 \\
R82 & +3.69 & +0.82 \\
I83 & +3.69 & +1.30 \\
K84 & +0.65 & +1.36 \\
H86 & +0.00 & +0.00 \\
Q87 & +0.00 & +0.00 \\
G88 & +2.76 & +2.67 \\
Q89 & +1.75 & +1.06 \\
H90 & +0.00 & +0.00 \\
I91 & +0.41 & +1.06 \\
G92 & +3.69 & +4.10 \\
E93 & +0.82 & +1.15 \\
M94 & +1.42 & +1.91 \\
\bottomrule
\end{tabular}
\end{table}


\subsection{MT-SP1 (matriptase)}
MT-SP1 / matriptase (ST14) is a type~II transmembrane serine protease of the S1 family; its
extracellular catalytic domain is over-expressed in epithelial cancers and is a validated
oncology target. The benchmark uses two inhibitory Fabs, E2 (\texttt{3BN9})~\cite{farady2008} and S4
(\texttt{3NPS})~\cite{schneider2012}, scanning the protease domain in chymotrypsin numbering
(Table~\ref{tab:ddg-mtsp1}; $\ddG$ via AB-Bind~\cite{sirin2016}). Simulation is
staged but \textbf{pending}: the crystal structures lack two internal residues (149, 218) that
must be repaired before MD (\S\ref{sec:repair}).
\begin{figure}[H]\centering\figorbox{mtsp1_structure.png}{0.45}
\caption{MT-SP1 protease domain (3NPS chain A). \textit{Placeholder.}}\end{figure}
% AUTO-GENERATED by analysis/make_label_tables.py -- do not edit by hand.
\begin{table}[htbp]\centering\small
\caption{\textbf{MT-SP1} alanine-scan \ddG\ (kcal/mol; enzymatic inhibition, $K_i$-derived via $\ddG=RT\ln(K^{\mathrm{mut}}/K^{\mathrm{wt}})$, $T=310$\,K; start structure: strip 3NPS chain A). Positive \ddG\ weakens binding.}
\label{tab:ddg-mtsp1}
\begin{tabular}{lrr}
\toprule
Residue & E2 (3BN9) & S4 (3NPS) \\
\midrule
Q23 & -- & +0.03 \\
I26 & -- & +0.65 \\
Q38 & -0.42 & -- \\
I41 & +0.00 & -- \\
I45 & -- & -0.34 \\
D46 & -- & +0.34 \\
D47 & -- & +1.08 \\
R48 & -- & -1.07 \\
F50 & -- & +0.22 \\
R51 & -- & +0.14 \\
Y52 & -- & +0.46 \\
I60 & +0.84 & -- \\
R82 & -- & -0.15 \\
R87 & -0.16 & -- \\
F89 & -- & +1.61 \\
N90 & -- & +0.26 \\
D91 & -- & +1.52 \\
F92 & -- & +0.47 \\
T93 & -- & +0.73 \\
F94 & +0.64 & -- \\
N95 & +0.77 & -- \\
D96 & +6.70 & -- \\
F97 & +6.70 & -- \\
T98 & +1.13 & -- \\
H138 & -- & +1.89 \\
Q140 & -- & +0.30 \\
Y141 & -- & +1.79 \\
H143 & +0.09 & -- \\
T144 & -- & +0.18 \\
Q145 & +0.13 & -- \\
Y146 & +1.09 & -- \\
L147 & -- & +0.30 \\
T150 & +0.29 & -- \\
L153 & +0.34 & -- \\
E163 & -- & +0.62 \\
Q168 & -- & -0.06 \\
E169 & +0.37 & +0.75 \\
Q174 & -0.03 & -- \\
Q175 & +2.51 & -- \\
D214 & -- & +1.48 \\
D217 & +0.57 & -- \\
Q218 & -- & -0.04 \\
R219 & -- & -0.08 \\
K221 & -- & -0.10 \\
R222 & -0.09 & -- \\
K224 & +0.79 & -- \\
\bottomrule
\end{tabular}
\end{table}


\subsection{Bont/A1 (H\textsubscript{C} receptor-binding domain)}
Botulinum neurotoxin serotype A1 is among the most potent toxins known: a $\sim$150\,kDa di-chain
protein comprising a zinc-metalloprotease light chain and a heavy chain. Its C-terminal
heavy-chain domain (H\textsubscript{C}, here residues 872--1295) is the receptor-binding domain
--- engaging neuronal gangliosides and the SV2 protein --- and the principal protective antigen
targeted by neutralizing antibodies. We simulate the isolated H\textsubscript{C} domain
(\texttt{2NYY} chain A); the catalytic light-chain zinc lies outside this fragment. Two
neutralizing mAbs, CR1 (\texttt{2NYY}) and AR2 (\texttt{2NZ9}), provide the labels
(Table~\ref{tab:ddg-bonta1})~\cite{garciarodriguez2007}; $\ddG$ via AB-Bind~\cite{sirin2016}.
\begin{figure}[H]\centering
\begin{minipage}[c]{0.37\linewidth}\centering\figorbox{bonta1_2NYY_ddg.png}{1.0}\\[-2pt]{\small CR1 (2NYY)}\end{minipage}\hfill
\begin{minipage}[c]{0.37\linewidth}\centering\figorbox{bonta1_2NZ9_ddg.png}{1.0}\\[-2pt]{\small AR2 (2NZ9)}\end{minipage}\hfill
\begin{minipage}[c]{0.13\linewidth}\centering\figorbox{ddgbar.png}{1.0}\end{minipage}
\caption{Bont/A1 H\textsubscript{C} coloured per residue by \ddG\ for two neutralizing mAbs (shared scale, white $=0$; grey $=$ not scanned). Red marks the hotspots --- the strongest reach \ddG\ $>7$ (binding essentially abolished).}\end{figure}
% AUTO-GENERATED by analysis/make_label_tables.py -- do not edit by hand.
\begin{table}[htbp]\centering\small
\caption{\textbf{Bont/A1} alanine-scan \ddG\ (kcal/mol; KinExA, $K_d$-derived via $\ddG=RT\ln(K^{\mathrm{mut}}/K^{\mathrm{wt}})$, $T=310$\,K; start structure: extract 2NYY Hc (res 872-1295)). Positive \ddG\ weakens binding.}
\label{tab:ddg-bonta1}
\begin{tabular}{lrr}
\toprule
Residue & mAb CR1 (2NYY) & mAb AR2 (2NZ9) \\
\midrule
S902 & -0.23 & -0.12 \\
K903 & +0.29 & +0.54 \\
Q915 & +0.52 & +0.10 \\
F917 & +0.42 & -0.05 \\
N918 & +0.89 & +2.16 \\
L919 & +2.59 & +2.28 \\
E920 & +2.84 & +2.77 \\
K923 & -0.14 & -0.30 \\
F953 & +4.06 & +3.34 \\
N954 & +0.10 & -0.15 \\
S955 & +0.00 & -0.08 \\
I956 & -0.01 & +0.07 \\
K1056 & -0.03 & -0.01 \\
D1058 & +0.01 & -0.03 \\
R1061 & +0.82 & +0.29 \\
D1062 & +2.53 & +2.34 \\
T1063 & +1.62 & +2.37 \\
H1064 & +7.32 & +7.42 \\
R1294 & +0.30 & +0.39 \\
\bottomrule
\end{tabular}
\end{table}


\subsection{IFN-$\gamma$ receptor}
The interferon-$\gamma$ receptor $\alpha$-chain (IFN-$\gamma$R1) ectodomain is the high-affinity
ligand-binding subunit of the type~II interferon receptor: two fibronectin type~III domains
forming an L-shaped $\beta$-sandwich that engages the IFN-$\gamma$ homodimer. The neutralizing
antibody A6 binds this ectodomain and the published alanine scan localises a compact energetic
hotspot on one face (Table~\ref{tab:ddg-ifngr})~\cite{hofstadter1999}; $\ddG$ values curated via
SKEMPI~2.0~\cite{skempi2}. It is a clean, rigid, single-chain
$\sim$95-residue Ig-like domain --- an attractive small test antigen. We simulate the apo
ectodomain stripped from \texttt{1JRH} (chain~I).
\begin{figure}[H]\centering
\begin{minipage}[c]{0.6\linewidth}\centering\figorbox{ifngr_1JRH_ddg.png}{0.95}\\[-2pt]{\small A6 (1JRH)}\end{minipage}\hfill
\begin{minipage}[c]{0.13\linewidth}\centering\figorbox{ddgbar.png}{1.0}\end{minipage}
\caption{IFN-$\gamma$ receptor coloured per residue by A6 alanine-scan \ddG\ (shared scale, white $=0$, red $=$ hotspot; grey $=$ not scanned). The hotspots cluster on a single $\beta$-sandwich face.}\end{figure}
% AUTO-GENERATED by analysis/make_label_tables.py -- do not edit by hand.
\begin{table}[htbp]\centering\small
\caption{\textbf{IFN-gamma receptor} alanine-scan \ddG\ (kcal/mol; SPR, $K_d$-derived via $\ddG=RT\ln(K^{\mathrm{mut}}/K^{\mathrm{wt}})$, $T=310$\,K; start structure: strip 1JRH chain I). Positive \ddG\ weakens binding.}
\label{tab:ddg-ifngr}
\begin{tabular}{lr}
\toprule
Residue & A6 (1JRH) \\
\midrule
K47 & +3.71 \\
N48 & +0.17 \\
Y49 & +3.53 \\
G50 & +4.45 \\
V51 & +1.68 \\
K52 & +3.39 \\
N53 & +4.30 \\
S54 & +0.38 \\
E55 & -0.57 \\
N79 & -0.20 \\
W82 & +4.43 \\
R84 & +0.39 \\
K98 & +0.31 \\
\bottomrule
\end{tabular}
\end{table}


\subsection{Tissue factor}
Tissue factor (TF, coagulation factor~III) is the cell-surface initiator of the extrinsic
blood-coagulation cascade: its extracellular region is a two-domain fibronectin type~III module
that binds and allosterically activates factor~VIIa. The inhibitory Fab 5G9 blocks the
macromolecular substrate site, and its antigen-side alanine scan marks the functional epitope
(Table~\ref{tab:ddg-tf})~\cite{huang1998}; $\ddG$ values curated via SKEMPI~2.0~\cite{skempi2}. The ectodomain is a well-folded, single-chain $\sim$200-residue
$\beta$-rich protein. We simulate the apo ectodomain stripped from \texttt{1AHW} (chain~C).
\begin{figure}[H]\centering
\begin{minipage}[c]{0.6\linewidth}\centering\figorbox{tf_1AHW_ddg.png}{0.95}\\[-2pt]{\small 5G9 (1AHW)}\end{minipage}\hfill
\begin{minipage}[c]{0.13\linewidth}\centering\figorbox{ddgbar.png}{1.0}\end{minipage}
\caption{Tissue factor coloured per residue by 5G9 alanine-scan \ddG\ (shared scale, white $=0$, red $=$ hotspot; grey $=$ not scanned).}\end{figure}
% AUTO-GENERATED by analysis/make_label_tables.py -- do not edit by hand.
\begin{table}[htbp]\centering\small
\caption{\textbf{Tissue factor} alanine-scan \ddG\ (kcal/mol; SPR, $K_d$-derived via $\ddG=RT\ln(K^{\mathrm{mut}}/K^{\mathrm{wt}})$, $T=310$\,K; start structure: strip 1AHW chain C). Positive \ddG\ weakens binding.}
\label{tab:ddg-tf}
\begin{tabular}{lr}
\toprule
Residue & 5G9 (1AHW) \\
\midrule
Y156 & +3.81 \\
Y157 & -1.89 \\
T167 & -0.07 \\
T170 & +1.11 \\
L176 & +0.99 \\
D178 & -0.48 \\
T197 & +1.35 \\
V198 & -0.31 \\
N199 & +1.08 \\
\bottomrule
\end{tabular}
\end{table}


\subsection{HCMV glycoprotein B}
Glycoprotein~B (gB) is the conserved class~III fusion machine of human cytomegalovirus and a
principal target of neutralizing antibodies. The ectodomain is a large, elongated trimer; the
antibody 1G2 recognises antigenic domain~2 and its alanine scan identifies the contact hotspots
(Table~\ref{tab:ddg-hcmvgb})~\cite{chandramouli2015}; $\ddG$ values curated via
SKEMPI~2.0~\cite{skempi2}. This is our largest antigen ($\sim$566 modelled residues) and is
biologically trimeric --- we simulate the apo protomer as deposited in \texttt{5C6T} (chain~A),
noting that the native oligomeric state may modulate the dynamics of the buried interfaces.
\begin{figure}[H]\centering
\begin{minipage}[c]{0.6\linewidth}\centering\figorbox{hcmvgb_5C6T_ddg.png}{0.95}\\[-2pt]{\small 1G2 (5C6T)}\end{minipage}\hfill
\begin{minipage}[c]{0.13\linewidth}\centering\figorbox{ddgbar.png}{1.0}\end{minipage}
\caption{HCMV glycoprotein~B protomer coloured per residue by 1G2 alanine-scan \ddG\ (shared scale, white $=0$, red $=$ hotspot; grey $=$ not scanned).}\end{figure}
% AUTO-GENERATED by analysis/make_label_tables.py -- do not edit by hand.
\begin{table}[htbp]\centering\small
\caption{\textbf{HCMV glycoprotein B} alanine-scan \ddG\ (kcal/mol; SPR, $K_d$-derived via $\ddG=RT\ln(K^{\mathrm{mut}}/K^{\mathrm{wt}})$, $T=310$\,K; start structure: strip 5C6T chain A). Positive \ddG\ weakens binding.}
\label{tab:ddg-hcmvgb}
\begin{tabular}{lr}
\toprule
Residue & 1G2 (5C6T) \\
\midrule
Y280 & +3.21 \\
N281 & +0.28 \\
T283 & +0.34 \\
N284 & +0.36 \\
N286 & -0.37 \\
F290 & +2.00 \\
E292 & +1.09 \\
F297 & +2.76 \\
F298 & +0.41 \\
I299 & +0.89 \\
\bottomrule
\end{tabular}
\end{table}


\subsection{Human growth hormone}
Human growth hormone (hGH, somatotropin) is a 191-residue four-helix-bundle cytokine and the
historical birthplace of alanine-scanning mutagenesis. Here the antibody-specific labels are the
MAb3 scan (Table~\ref{tab:ddg-hgh})~\cite{jin1992}\,\textbf{[verify]} (benchmark cites L.\ Jin's
1994 thesis, Table~3.2; the published values are in~\cite{jin1992}); the benchmark additionally
carries a secondary,
provisional profile of the total \ddG\ aggregated across 21 anti-hGH MAbs (not antibody-specific,
held out of the primary labels). We simulate the apo hormone from \texttt{1HGU} (chain~A).
\begin{figure}[H]\centering
\begin{minipage}[c]{0.6\linewidth}\centering\figorbox{hgh_1HGU_ddg.png}{0.95}\\[-2pt]{\small MAb3 (1HGU)}\end{minipage}\hfill
\begin{minipage}[c]{0.13\linewidth}\centering\figorbox{ddgbar.png}{1.0}\end{minipage}
\caption{Human growth hormone coloured per residue by MAb3 alanine-scan \ddG\ (shared scale, white $=0$, red $=$ hotspot; grey $=$ not scanned). Only the four MAb3-scanned positions are coloured.}\end{figure}
% AUTO-GENERATED by analysis/make_label_tables.py -- do not edit by hand.
\begin{table}[htbp]\centering\small
\caption{\textbf{human growth hormone} alanine-scan \ddG\ (kcal/mol; ELISA, $K_d$-derived via $\ddG=RT\ln(K^{\mathrm{mut}}/K^{\mathrm{wt}})$, $T=310$\,K; start structure: apo 1HGU). Positive \ddG\ weakens binding.}
\label{tab:ddg-hgh}
\begin{tabular}{lr}
\toprule
Residue & MAb3 (1HGU) \\
\midrule
R8 & +2.16 \\
R16 & +3.68 \\
D112 & +1.84 \\
D116 & +2.73 \\
\bottomrule
\end{tabular}
\end{table}

